% ===========================================================================
% Poste 6 — Frais de garde (crédit QC + déduction fédérale + coût)
% ===========================================================================
\poste{6}{Frais de garde d'enfants (crédit, déduction et coût)}{QC\_garde, Frais\_garde}

\pdesc Les frais de garde payés par le ménage ont \emph{trois} effets sur le
revenu disponible : un crédit d'impôt remboursable au Québec (sur les frais
\emph{non subventionnés} seulement), une déduction au fédéral (sur tous les
frais) et, bien sûr, le coût lui-même, soustrait du revenu disponible.

\pobj Soutenir les parents qui travaillent ou étudient en remboursant une
part substantielle des frais de garde --- le crédit québécois complète la
place subventionnée à contribution réduite en couvrant la garde non
subventionnée.

\paragraph{Règle de calcul --- 1) crédit d'impôt remboursable du Québec.}
Le taux du crédit
$\tau(\RFN)$ décroît par paliers de \pc{78} à \pc{67} selon le revenu
familial net (tableau ci-dessous). Les frais admissibles sont les frais
\emph{non subventionnés}, bornés par un plafond qui dépend de l'âge de
l'enfant ; le plafonnement est \emph{agrégé} sur la famille :
\[
\text{crédit} = \tau(\RFN) \times
\min\Bigl(\sum_{e \,:\, f_e > 0} p(\text{âge}_e),\;\; \sum_e f_e\Bigr)
\]
où $f_e$ est le total des frais non subventionnés de l'enfant $e$ et
$p(\text{âge})$ son plafond annuel : le plafond \og jeune \fg{}
(\D{12275} en 2025, \D{12525} en 2026) si l'âge est de 5~ans ou moins, le
plafond \og autre enfant admissible \fg{} (\D{6180} / \D{6305}) jusqu'à
l'âge d'admissibilité, zéro au-delà. L'âge limite passe de 16~ans (2025) à
14~ans (2026), resserrement annoncé du programme.

\paragraph{Règle de calcul --- 2) déduction fédérale pour frais de garde.}
Tous les frais payés (subventionnés ou non) sont déductibles au fédéral,
dans une triple limite :
\[
D_{garde} = \min\Bigl(\sum_e f_e^{tot},\;\;
\sum_{e \,:\, f^{tot}_e > 0} p^{\text{féd}}(\text{âge}_e),\;\;
\tfrac{2}{3} \times r_{\min}\Bigr)
\]
où $p^{\text{féd}}$ vaut \D{8000} (âge $\le 5$), \D{5000} (6 à 15~ans) ou zéro
(16~ans et plus), et $r_{\min}$ est le revenu de travail du conjoint le
moins payé (celui qui doit réclamer la déduction). $D_{garde}$ réduit le
revenu imposable fédéral \emph{et} l'$\AFNI$
(équation~\ref{eq:afni}) --- il augmente donc l'Allocation canadienne pour
enfants et le crédit TPS en plus de réduire l'impôt fédéral.

\paragraph{Règle de calcul --- 3) coût.} Le total des frais payés,
subventionnés et non subventionnés, est soustrait du revenu disponible
(terme $G$ de l'équation~\ref{eq:rd}).

\pparams Barème du taux du crédit québécois selon le revenu familial net :

\begin{center}
\small
\begin{tabular}{r r c}
\toprule
$\RFN \le$ (2025) & $\RFN \le$ (2026) & Taux \\
\midrule
\D{24795}  & \D{25305}  & \pc{78} \\
\D{43725}  & \D{44620}  & \pc{75} \\
\D{45340}  & \D{46270}  & \pc{74} \\
\D{46970}  & \D{47935}  & \pc{73} \\
\D{48570}  & \D{49565}  & \pc{72} \\
\D{50195}  & \D{51225}  & \pc{71} \\
\D{119835} & \D{122290} & \pc{70} \\
au-delà    & au-delà    & \pc{67} \\
\bottomrule
\end{tabular}
\hspace{2em}
\begin{tabular}{l r r}
\toprule
Plafonds & 2025 & 2026 \\
\midrule
Crédit QC --- enfant de 5 ans ou moins & \D{12275} & \D{12525} \\
Crédit QC --- autre enfant admissible  & \D{6180}  & \D{6305} \\
Crédit QC --- âge limite (exclusif)    & 16 ans    & 14 ans \\
Déduction féd. --- âge $\le 5$         & \D{8000}  & \D{8000} \\
Déduction féd. --- 6 à 15 ans          & \D{5000}  & \D{5000} \\
Déduction féd. --- borne de revenu     & $2/3 \times r_{\min}$ & $2/3 \times r_{\min}$ \\
\bottomrule
\end{tabular}
\end{center}

\pnotes Deux choix du calculateur officiel sont reproduits par fidélité :
(1)~le seuil \og jeune \fg{} est l'âge $\le 5$, alors que la règle légale
vise l'enfant \og de moins de 7~ans \fg{} (âge $\le 6$) --- divergence d'un
an ; (2)~le plafonnement est \emph{agrégé} (somme des plafonds contre somme
des frais) plutôt qu'enfant par enfant, ce qui peut surévaluer les frais
admissibles quand un enfant dépasse son plafond et qu'un autre est en deçà.
Le plafond majoré pour enfant handicapé n'est pas modélisé.

% ============================================================================
% GÉNÉRÉ par scripts/exemples-guide.ts — NE PAS ÉDITER À LA MAIN.
% Régénérer : npm run exemples (chaque chiffre est vérifié contre le moteur).
% ============================================================================
\pexemples Trois effets distincts sur le revenu disponible : le \textbf{crédit
québécois} (frais \emph{non subventionnés} seulement), la \textbf{déduction
fédérale} (tous les frais ; elle réduit l'$\AFNI$, donc augmente l'ACE et le
crédit TPS, et réduit l'impôt fédéral — postes~14, 15 et 19), et le
\textbf{coût} des frais payés, soustrait du revenu disponible.

\medskip\noindent\textbf{M6 --- famille monoparentale, 35~ans, revenu de travail de \D{35000}, un enfant de 3~ans en garderie subventionnée (\D{2000} payés dans l'année).}
Garde \emph{subventionnée} : aucun frais admissible au crédit
québécois (crédit nul). La déduction fédérale, elle, s'applique :
\begin{align*}
\text{déduction} &= \min\bigl(\num{2000},\ \num{8000},\ \tfrac{2}{3} \times \num{35000}\bigr) = \num{2000}
\end{align*}
Elle abaisse l'$\AFNI$ de \num{34685} à \num{32685} (le poste~19
détaille la soustraction de la cotisation RRQ supplémentaire). Enfin, les
\D{2000} payés sont soustraits du revenu disponible.

\medskip\noindent\textbf{M8 --- couple, 38 et 36~ans, revenus de travail de \D{60000} et \D{40000}, deux enfants : 4~ans (garde non subventionnée, \D{13000}) et 8~ans (garde non subventionnée, \D{3000}).}
Frais non subventionnés de \D{13000} (4~ans) et \D{3000}
(8~ans). Le plafonnement québécois est \emph{agrégé} : on additionne d'abord
les plafonds des enfants concernés, puis on borne la \emph{somme} des frais.
\begin{align*}
\text{plafonds} &= \num{12275} + \num{6180} = \num{18455} \\
\text{frais admissibles} &= \min(\num{18455},\ \num{16000}) = \num{16000} \\
\text{crédit} &= \pc{70} \times \num{16000} = \num{11200}
\end{align*}
Le taux de \pc{70} découle du $\RFN$ de \num{96230}
(barème de 78\,\% à 67\,\%). Noter l'effet du plafond agrégé : l'enfant de
4~ans dépasse à lui seul son plafond de \num{12275}, mais l'excédent
est « absorbé » par le plafond inutilisé de l'autre enfant — les \num{16000}
sont admissibles en entier. Au fédéral, le plafonnement est agrégé aussi :
\begin{align*}
\text{déduction} &= \min\bigl(\num{16000},\ \num{8000} + \num{5000},\ \tfrac{2}{3} \times \num{40000}\bigr) = \num{13000}
\end{align*}
(la borne de $\tfrac{2}{3}$ porte sur le revenu de travail du conjoint le
moins payé). Coût soustrait du revenu disponible : \D{16000}.

\medskip\noindent\textbf{M9 --- couple, 45 et 45~ans, revenus de travail de \D{120000} et \D{60000}, un enfant de 5~ans (garde non subventionnée, \D{15000}).}
$\RFN$ de \num{175521} : au-delà du dernier seuil du
barème, le taux est au \emph{plancher} de \pc{67}.
\begin{align*}
\text{frais admissibles} &= \min(\num{12275},\ \num{15000}) = \num{12275} \\
\text{crédit} &= \pc{67} \times \num{12275} = \num{8224.25}
\end{align*}
Déduction fédérale : $\min(\num{15000},\ \num{8000},\
\tfrac{2}{3} \times \num{60000}) = \num{8000}$. Coût soustrait : \D{15000}.


\prefs Loi sur les impôts (RLRQ, c.~I-3), art.~1029.8.67 et suivants ;
déclaration TP-1, annexe~C ; Revenu Québec. Fédéral : Loi de l'impôt sur le
revenu (L.R.C. 1985, ch.~1 (5\ieme{}~suppl.)), art.~63 ; formulaire T778 ;
ARC.
